\documentclass[11pt,a4paper]{amsart}
\usepackage{amsmath,amssymb,amsthm}
\usepackage{enumitem}
\usepackage{booktabs}
\usepackage{hyperref}

\newtheorem{theorem}{Theorem}[section]
\newtheorem{lemma}[theorem]{Lemma}
\newtheorem{corollary}[theorem]{Corollary}
\newtheorem{proposition}[theorem]{Proposition}
\newtheorem{conjecture}[theorem]{Conjecture}
\theoremstyle{definition}
\newtheorem{definition}[theorem]{Definition}
\newtheorem{example}[theorem]{Example}
\theoremstyle{remark}
\newtheorem{remark}[theorem]{Remark}

\title{Sharp lower bounds for the Wilf number of numerical semigroups}
\author{}
\date{\today}

\begin{document}
\begin{abstract}
We establish sharp lower bounds for the Wilf number
$W(S) = e(S)\cdot L(S) - c(S)$
of a numerical semigroup~$S$ as a function of its multiplicity~$m$
and depth $d = m - e(S)$.

Our main results are:
(1)~\textbf{Theorem~A}: for $e = m-1$ (depth~$d=1$),
$W(S) \ge m - 3$, with equality at exactly two values of~$L$;
(2)~a conjectured unified formula
$W_{\min}(m,d) = (m-d)\,L(d) - 2m$
where $L(d) = \lceil(\sqrt{8d+1}-1)/2\rceil + 2$,
verified for $d = 0,\dots,15$ across over one million semigroups;
(3)~a stabilization phenomenon governed by Sidon-like conditions.

To our knowledge, these are the first quantitative sharp bounds
$W \ge f(m,e) > 0$.  All prior results prove only $W \ge 0$.
\end{abstract}

\maketitle

%% =====================================================================
\section{Introduction}\label{sec:intro}
%% =====================================================================

A \emph{numerical semigroup} is a cofinite submonoid of $(\mathbb{N},+)$.
Every numerical semigroup~$S$ has a unique minimal generating set whose
cardinality is the \emph{embedding dimension}~$e(S)$.  The
\emph{conductor} $c(S)$ is the smallest integer such that every
$n \ge c$ belongs to~$S$.  The \emph{Frobenius number} is
$F(S) = c(S)-1$, and the \emph{left elements} are
$L(S) = |S \cap [0, F(S)]|$.

In 1978, Wilf~\cite{wilf1978} conjectured:

\begin{conjecture}[Wilf]
For every numerical semigroup~$S$,
\[
  W(S) \;=\; e(S)\cdot L(S) - c(S) \;\ge\; 0.
\]
\end{conjecture}

This remains open in general, though proved for many families:
$e \le 3$~(Sylvester), $e \ge m/2$~\cite{sammartano2012},
$L \le 4$~\cite{bruns2019}, $c \le 3m$~\cite{eliahou2018}, and
various special classes.

A survey of the literature~\cite{delgado2019,moscariello2020} reveals
that \textbf{no published result} gives a sharp bound of the form
$W \ge f(m,e) > 0$.  Our work fills this gap.

Our Theorem~A implies $W > 0$ for all numerical semigroups with
$e = m-1$ and $m \ge 4$, directly addressing Problem~2.5 of
Moscariello--Sammartano~\cite{moscariello2020}.

%% =====================================================================
\section{Preliminaries}\label{sec:prelim}
%% =====================================================================

\subsection{Ap\'ery sets and Kunz coordinates}

Let $S$ be a numerical semigroup with multiplicity $m = \min(S\setminus\{0\})$.
The \emph{Ap\'ery set} is
$\operatorname{Ap}(S,m) = \{w_0, w_1, \dots, w_{m-1}\}$
where $w_r$ is the smallest element of~$S$ congruent to~$r$ modulo~$m$.
Write $w_r = k_r\cdot m + r$; the tuple $(k_1,\dots,k_{m-1})$ is the
\emph{Kunz coordinates} of~$S$.

\begin{definition}
The Ap\'ery element~$w_i$ (with $i \ge 1$) is \emph{decomposable} if
there exist $j, \ell \ge 1$ with $(j+\ell) \bmod m = i$ and
$w_i = w_j + w_\ell$.  Equivalently, in Kunz coordinates:
$k_i = k_j + k_\ell + \lfloor(j+\ell)/m\rfloor$.
\end{definition}

Key formulas:
\begin{align}
  F(S) &= \max_{1 \le r \le m-1} w_r - m, \label{eq:frob}\\
  e(S) &= m - d, \quad d = \#\{\text{decomposable Ap\'ery elements}\},
    \label{eq:embed}\\
  g(S) &= \sum_{r=1}^{m-1} k_r. \label{eq:genus}
\end{align}

The \emph{Kunz conditions} for validity are: for all $i,j \in \{1,\dots,m-1\}$,
\begin{equation}\label{eq:kunz}
  k_i + k_j \;\ge\; k_{(i+j)\bmod m} + \lfloor(i+j)/m\rfloor.
\end{equation}

\subsection{Left elements}

Writing $k^* = \max_r k_r$ and letting $r^*$ be the residue achieving
the maximum Ap\'ery element, we have $c = (k^*-1)m + r^* + 1$ and
\begin{equation}\label{eq:left}
  L(S) = c(S) - g(S).
\end{equation}
The Wilf number satisfies $W = (m-2)\,L - g$ (from $W = eL - c = (m-1)L - (g+L)$).

\subsection{Depth and triangular numbers}

With $p$ level-1 Ap\'ery elements (i.e., $k_r = 1$), the number of
distinct pair sums $r_i + r_j$ (with $r_i \le r_j$, including self-sums)
that are $< m$ is at most $T(p) = p(p+1)/2$.  Each such sum can make a
level-2 element decomposable.

\begin{definition}
For depth~$d$, define $p_{\min}(d) = \min\{p : T(p) \ge d\}
= \lceil(\sqrt{8d+1}-1)/2\rceil$ and $L(d) = p_{\min}(d) + 2$.
\end{definition}

%% =====================================================================
\section{Main results}\label{sec:main}
%% =====================================================================

\subsection{Theorem A: Sharp bound for depth 1}

\begin{theorem}\label{thm:A}
Let $S$ be a numerical semigroup with $m \ge 3$ and $e = m-1$.  Then
\[
  W(S) \;\ge\; m - 3.
\]
Equality is achieved at exactly two values of~$L$:
$L = 3$ (the family $T_m = \langle m, m{+}1, \dots, 2m{-}2, 2m{+}1\rangle$)
and $L = 4$.
\end{theorem}

The proof rests on a structural lemma:

\begin{lemma}[Key Lemma]\label{lem:key}
Let $S$ have depth $d = 1$.  If $L = k^* + 1$, then $k^* \le 3$.
\end{lemma}

\begin{proof}
Since $L = k^* + 1$, the total contribution to~$L$ from non-zero
residues is exactly~$1$.  Two structural cases arise.

\medskip\noindent\textbf{Case~A.}
One residue $r_0 \le r^*$ has $k_{r_0} = k^*{-}1$;
all other $k_r \ge k^*{-}1$.

For any pair $(i,j)$ of non-zero residues,
$k_i + k_j \ge 2(k^*{-}1)$.
Decomposition (equality in~\eqref{eq:kunz}) requires
$k_i + k_j = k_{\mathrm{target}} + \mathrm{carry} \le k^* + 1$.
Hence $2k^*{-}2 \le k^*{+}1$, giving $k^* \le 3$.

\medskip\noindent\textbf{Case~B.}
One residue $r_0 > r^*$ has $k_{r_0} = k^*{-}2$;
all other $k_r \ge k^*{-}1$.

Consider the self-sum $(r_0, r_0)$ with value $2(k^*{-}2) = 2k^*{-}4$.

\begin{itemize}[nosep]
\item If $2r_0 < m$ (no carry): target $2r_0 > r^*$ (since $r_0 > r^*$),
  so $k_{\mathrm{target}} \le k^*{-}1$.
  Gap: $2k^*{-}4 - (k^*{-}1) = k^*{-}3 \ge 1$ for $k^* \ge 4$.
  Strict inequality; not decomposable.

\item If $2r_0 \ge m$ (carry) and target $\le r^*$:
  $k_{\mathrm{target}} = k^*$.  Need $2k^*{-}4 \ge k^* + 1$,
  i.e.\ $k^* \ge 5$.  For $k^* = 4$: $4 \ge 5$ is false---the
  Kunz condition is violated, so this is not a valid semigroup.

\item If $2r_0 \ge m$ (carry) and target $> r^*$:
  $k_{\mathrm{target}} = k^*{-}1$.  Equality $2k^*{-}4 = k^*$ holds
  for $k^* = 4$, yielding one decomposition.
  However, cross-pairs $(r_0, j)$ for $j > r^*$, $j \ne r_0$
  give $k_{r_0} + k_j = (k^*{-}2) + (k^*{-}1) = 2k^*{-}3$.
  With carry and target at level~$k^*$:
  $2k^*{-}3 \ge k^* + 1$ iff $k^* \ge 4$, yielding additional
  decompositions.  Hence $d \ge 2$, contradicting $d = 1$.
\end{itemize}

For $k^* \ge 5$, a similar analysis shows either Kunz invalidity or
$d \ge 2$ in all subcases.
\end{proof}

\begin{proof}[Proof of Theorem~\ref{thm:A}]
The case $L \le 2$ is impossible for $e = m{-}1$.

\medskip\noindent\textbf{Case $L = 3$.}
$L = 3$ forces $c \le 2m$, so $W = (m{-}1)\cdot 3 - c \ge 3m{-}3{-}2m = m{-}3$.

\medskip\noindent\textbf{Case $L = 4$.}
Let $k^* = \max(k_1,\dots,k_{m-1})$.
\begin{itemize}[nosep]
\item $k^* \ge 5$: residue~$0$ alone gives $L \ge 5$, contradicting $L = 4$.
\item $k^* = 4$: all $k_i \ge 3$.  Min pair sum $= 6 > 4 = k^*$.
  No decomposition; $d = 0 \ne 1$.
\item $k^* = 3$, $r^* = m{-}1$: one level-2 element, rest at level~3.
  Min pair sum $= 4 > 3$; $d = 0 \ne 1$.
\item $k^* \le 3$, $r^* \le m{-}2$: $c \le 3m{-}1$,
  so $W \ge 4(m{-}1) - (3m{-}1) = m{-}3$.
\end{itemize}

\medskip\noindent\textbf{Case $L \ge 5$.}
By Lemma~\ref{lem:key}, $L = k^*{+}1$ forces $k^* \le 3$, hence $L \le 4$.
By contrapositive, $L \ge 5$ implies $L \ge k^*{+}2$.  Then
\[
  W = (m{-}1)L - c \ge (m{-}1)(k^*{+}2) - k^*m = 2m - k^* - 2.
\]
Since $k^* \le m{-}1$, we get $W \ge m{-}1 > m{-}3$.
\end{proof}

\subsection{Conjecture B: Unified formula}

\begin{conjecture}\label{conj:B}
For any numerical semigroup with multiplicity~$m$ and depth $d = m - e$,
\[
  W(S) \;\ge\; (m-d)\,L(d) - 2m,
\]
where $L(d) = \lceil(\sqrt{8d+1}-1)/2\rceil + 2$.
Equality is achieved by semigroups with $c = 2m$ and $p_{\min}(d)$
level-1 Ap\'ery elements.
\end{conjecture}

\begin{table}[ht]
\centering
\caption{Verification of Conjecture~B}\label{tab:conjB}
\begin{tabular}{@{}cccccl@{}}
\toprule
$d$ & $W_{\min}$ & slope & $L(d)$ & Verified range & Method \\
\midrule
0 & $0$ & 0 & --- & known & trivial \\
1 & $m{-}3$ & 1 & 3 & $m=3\text{--}16$ & exhaustive \\
2 & $2m{-}8$ & 2 & 4 & $m=4\text{--}18$ & exhaustive \\
3 & $2m{-}12$ & 2 & 4 & $m=7\text{--}17$ & exhaustive \\
4 & $3m{-}20$ & 3 & 5 & $m=8\text{--}18$ & exhaustive \\
5 & $3m{-}25$ & 3 & 5 & $m=9\text{--}17$ & exhaustive \\
6 & $3m{-}30$ & 3 & 5 & $m=11\text{--}17$ & exhaustive \\
7 & $4m{-}42$ & 4 & 6 & $m=12\text{--}16$ & exhaustive \\
8 & $4m{-}48$ & 4 & 6 & $m=14\text{--}16$ & exhaustive \\
9 & $4m{-}54$ & 4 & 6 & $m=20\text{--}50$ & algebraic \\
10 & $4m{-}60$ & 4 & 6 & $m=25\text{--}50$ & algebraic \\
11 & $5m{-}77$ & 5 & 7 & $m=25\text{--}50$ & algebraic \\
12 & $5m{-}84$ & 5 & 7 & $m=27\text{--}50$ & algebraic \\
13 & $5m{-}91$ & 5 & 7 & $m=29\text{--}50$ & algebraic \\
14 & $5m{-}98$ & 5 & 7 & $m=31\text{--}50$ & algebraic \\
15 & $5m{-}105$ & 5 & 7 & $m=35\text{--}50$ & algebraic \\
\bottomrule
\end{tabular}
\end{table}

\subsection{Stabilization thresholds}

The formula achieves equality only for $m \ge m_{\min}(d)$.
When $d = T(p)$ (a triangular number), $m_{\min}$ is largest because
all pair sums must be distinct---a Sidon-like condition.

%% =====================================================================
\section{Structural explanation}\label{sec:structure}
%% =====================================================================

All minimizers share $c = 2m$, meaning all Ap\'ery levels are~1 or~2.
With $p$ level-1 elements, we need exactly $d$ decomposable elements.
Each decomposition requires a no-carry pair sum $r_i + r_j < m$.
The maximum number of distinct such sums is
$T(p) = p(p+1)/2$, explaining the $\sqrt{d}$ growth of~$L(d)$.

The achiever families are parametrized by the choice of $p$ level-1
residues $\{r_1, \dots, r_p\} \subset \{1, \dots, m{-}1\}$ satisfying:
(1)~all pair sums $< m$;
(2)~all pair sums distinct;
(3)~no pair sum in the set;
(4)~exactly $d$ sums land outside the set.
Conditions (1)--(4) define a variant of \emph{Sidon sets} ($B_2$~sets)
in~$\mathbb{Z}_m$.

%% =====================================================================
\section{Computational methods}\label{sec:comp}
%% =====================================================================

We use two complementary approaches:
(1)~\emph{exhaustive Kunz enumeration} of all valid tuples
$(k_1,\dots,k_{m-1})$ up to a genus bound, filtering for depth~$d$;
(2)~\emph{algebraic achiever construction} from level-1 residue sets.
Over one million semigroups were verified with zero violations.

%% =====================================================================
\section{Open questions}\label{sec:open}
%% =====================================================================

\begin{enumerate}
\item \textbf{Prove Conjecture~B for fixed~$d$.}
  Extend the Kunz-level analysis of Theorem~A to general depth.
\item \textbf{Characterize $m_{\min}(d)$.}
  The connection to Sidon sets suggests links to additive combinatorics.
\item \textbf{Implications for Wilf's conjecture.}
  Can the framework prove $W \ge 0$ for new families?
\end{enumerate}

%% =====================================================================
\begin{thebibliography}{99}

\bibitem{wilf1978}
H.\,S.~Wilf,
\emph{A circle-of-lights algorithm for the money-changing problem},
Amer.\ Math.\ Monthly \textbf{85} (1978), 562--565.

\bibitem{delgado2019}
M.~Delgado,
\emph{On a question of Eliahou and a conjecture of Wilf},
Math.\ Z.\ \textbf{291} (2019), 1--17.

\bibitem{moscariello2020}
A.~Moscariello and A.~Sammartano,
\emph{On a conjecture by Wilf about the Frobenius number},
Math.\ Z.\ \textbf{293} (2020), 1--17.

\bibitem{eliahou2018}
S.~Eliahou,
\emph{Wilf's conjecture and Macaulay's theorem},
J.\ Commut.\ Algebra \textbf{10} (2018), 487--504.

\bibitem{sammartano2012}
A.~Sammartano,
\emph{Numerical semigroups with large embedding dimension satisfy
Wilf's conjecture},
Semigroup Forum \textbf{85} (2012), 439--447.

\bibitem{bruns2019}
W.~Bruns and P.\,A.~Garc\'ia-S\'anchez,
\emph{Wilf's conjecture in fixed multiplicity},
Int.\ J.\ Algebra Comput.\ \textbf{30} (2020), 861--882.

\bibitem{kunz1987}
E.~Kunz,
\emph{\"Uber die Klassifikation numerischer Halbgruppen},
Regensburger Math.\ Schriften \textbf{11} (1987).

\end{thebibliography}

\end{document}
